% Presentación de primera clase — Máquinas térmicas para ingenieros
% Pontificia Universidad Javeriana — Prof. Camilo Bayona
% Compilar con: pdflatex (dos pasadas) — no requiere paquetes exóticos.
\documentclass[aspectratio=169]{beamer}

\usetheme{Madrid}
\usecolortheme{beaver}
\usefonttheme{professionalfonts}

\usepackage[utf8]{inputenc}
\usepackage[T1]{fontenc}
\usepackage[spanish,es-noindentfirst]{babel}
\usepackage{booktabs}
\usepackage{tabularx}
\usepackage{amsmath}
\usepackage{array}

\setbeamertemplate{navigation symbols}{}
\setbeamertemplate{footline}[frame number]

\title[Máquinas Térmicas — Rediseño 2026]{Máquinas térmicas para ingenieros}
\subtitle{Rediseño metodológico: AVA + cátedra invertida + proyecto aplicado}
\author{Camilo Bayona}
\institute[PUJ]{Pontificia Universidad Javeriana \\ Facultad de Ingeniería — Departamento de Ingeniería Industrial}
\date{Primera sesión, 2026}

\begin{document}

% ---------------------------------------------------------------
\begin{frame}
  \titlepage
\end{frame}

% ---------------------------------------------------------------
\begin{frame}{Agenda de hoy}
  \tableofcontents
\end{frame}

% =================================================================
\section{Contexto}
% =================================================================
\begin{frame}{¿Por qué actualizamos el curso?}
  \begin{itemize}
    \item El curso ya cubre ciclos reales, combustión, calderas, motores reciprocantes, turbinas de gas,
      plantas de vapor, cogeneración y refrigeración.
    \item Se conserva \textbf{todo} el alcance disciplinar aprobado (ID CURSO 036053).
    \item Lo que cambia es \textbf{cómo se entrega, se discute y se aplica} ese contenido.
    \item Cuatro decisiones pedagógicas, con evaluación cuantitativa (rigurosidad de ingeniería) y
      cualitativa (liderazgo de proyecto y de negocio en equipo).
  \end{itemize}
\end{frame}

\begin{frame}{Modelo: Cátedra — Tecnología — Proyecto}
  \begin{columns}[T]
    \column{0.48\textwidth}
      \textbf{Antes}
      \begin{itemize}
        \item Clase magistral en el aula.
        \item Notas de clase estáticas.
        \item Evaluación centrada en parciales.
      \end{itemize}
    \column{0.48\textwidth}
      \textbf{Ahora}
      \begin{itemize}
        \item OVA interactivos (Jupyter + CoolProp) en Teams desde el día 1.
        \item Aula = discusión invertida.
        \item Proyecto aplicado: maqueta de motor + modelo de negocio.
      \end{itemize}
  \end{columns}
\end{frame}

% =================================================================
\section{Cuatro ideas pedagógicas}
% =================================================================
\begin{frame}{Idea 1 — AVA en Teams desde el primer día}
  \begin{itemize}
    \item Todos los OVA (cuadernos + datos + bibliografía) se publican en Teams \textbf{desde la sesión 1},
      no semana a semana.
    \item El estudiante siempre tiene el cuadernillo disponible para leerlo y prepararlo \textbf{antes} de
      cada sesión presencial.
    \item Permite avanzar a ritmo propio y llegar a clase con preguntas concretas.
  \end{itemize}
\end{frame}

\begin{frame}{Idea 2 — Sesión presencial invertida}
  \begin{itemize}
    \item El aula \textbf{no repite} lo que ya está en el OVA.
    \item Se proyectan las temáticas y se discuten activamente entre el profesor y los estudiantes.
    \item Objetivo: la sesión más invertida, participativa y \textbf{divertida} posible.
    \item El tiempo de aula se usa para resolver dudas, profundizar casos y conectar con el proyecto.
  \end{itemize}
\end{frame}

\begin{frame}{Idea 3 — Proyecto aplicado: maqueta didáctica de un motor térmico}
  \begin{enumerate}
    \item \textbf{Memoria de cálculo}: modelación termodinámica del motor en Jupyter + CoolProp,
      extendiendo los OVA del curso.
    \item \textbf{Diseño geométrico}: modelador CAD y ensamblador de proyectos mecánicos, trazable a la
      memoria de cálculo.
    \item \textbf{Fabricación y ensamblaje}: impresión 3D en PETG (material del profesor) y ensamblaje en
      las máquinas herramientas del \textbf{CTAI}.
  \end{enumerate}
\end{frame}

\begin{frame}{Idea 4 — Modelo de negocio}
  \begin{itemize}
    \item Cada equipo concibe un \textbf{modelo de negocio} basado en las capacidades de ingeniería térmica
      del curso.
    \item Debe demostrar \textbf{pertinencia y viabilidad comercial en el mercado nacional}.
    \item Se sustenta junto con la maqueta y su memoria de cálculo: la solución técnica se defiende también
      como propuesta de valor.
  \end{itemize}
\end{frame}

% =================================================================
\section{Evaluación}
% =================================================================
\begin{frame}{Evaluación: cuantitativa + cualitativa}
  \begin{table}
    \small
    \begin{tabularx}{\linewidth}{lcXX}
      \toprule
      \textbf{Corte} & \textbf{Peso} & \textbf{Cuantitativo (rigor técnico)} & \textbf{Cualitativo (liderazgo)} \\
      \midrule
      1 & 30\% & Talleres/parcial ciclos + combustión; memoria de cálculo preliminar & Constitución de
        equipo y cronograma; selección del negocio. \\
      2 & 30\% & Parcial-taller calderas + reciprocantes; diseño CAD trazable & Avance de fabricación en
        CTAI; gestión de equipo; avance del negocio. \\
      3 & 40\% & Maqueta fabricada + memoria completa; examen final & Sustentación del modelo de negocio;
        liderazgo del proyecto \\
      \bottomrule
    \end{tabularx}
  \end{table}
\end{frame}

% =================================================================
\section{Los OVA del curso}
% =================================================================
\begin{frame}{OVA ya construidos y verificados (bloques 1–5)}
  \small
  \begin{columns}[T]
    \column{0.5\textwidth}
      \begin{enumerate}
        \item Procesos Politrópicos
        \item Ciclos Termodinámicos
        \item Ciclos Biogeoquímicos
        \item Combustibles Fósiles
        \item Repaso de Química
        \item[5b.] Reacciones Nucleares \emph{(frontera)}
        \item Cinética de Reacción
      \end{enumerate}
    \column{0.5\textwidth}
      \begin{enumerate}
        \setcounter{enumi}{7}
        \item Modelación de la Combustión
        \item Cinética de Combustión
        \item[8b.] Combustión de Biomasa \emph{(frontera)}
        \item Tecnología de la Combustión
        \item Compresores y Turbinas Reciprocantes
        \item Motores Reciprocantes
        \item Ensayo de Motores
      \end{enumerate}
  \end{columns}
  \vspace{0.3cm}
  \footnotesize Cada uno ejecuta sin errores y verifica con error $<1\%$ frente a un caso de referencia.
\end{frame}

\begin{frame}{Lo que sigue viéndose cada semestre — Bloque 6}
  \textbf{Motores de turbina de gas} \emph{(OVA en construcción; por ahora, cátedra)}
  \begin{itemize}
    \item Ciclo \textbf{Brayton}: ideal, regenerativo, con recalentamiento, con enfriamiento intermedio en
      la compresión.
    \item Compresor centrífugo y axial; cámara de combustión; turbina radial y axial.
    \item Propulsión aeronáutica: turborreactor, turboventilador, turbohélice, Ramjet, Scramjet.
    \item Motores de autopropulsión: ecuación fundamental de la astronáutica, tobera, motor químico
      (sólido, líquido, híbrido).
  \end{itemize}
\end{frame}

\begin{frame}{Lo que sigue viéndose cada semestre — Bloques 7 a 9}
  \small
  \textbf{Bloque 7 — Plantas térmicas con turbina de vapor}
  \begin{itemize}
    \item Ciclo \textbf{Rankine}: ideal/real, con sobrecalentamiento, con recalentamiento, con
      regeneración, supercrítico. Turbinas de vapor, calderas, condensador y bomba de presión.
  \end{itemize}
  \textbf{Bloque 8 — Plantas de cogeneración}
  \begin{itemize}
    \item Tipologías (turbina de gas, vapor, reciprocante, ciclo combinado); rendimientos, ahorro de
      combustible.
  \end{itemize}
  \textbf{Bloque 9 — Máquinas de refrigeración y aire acondicionado}
  \begin{itemize}
    \item Ciclo de compresión de vapor (\textbf{Rankine inverso}), absorción, refrigeración evaporativa;
      psicrometría, temperatura de bulbo seco/húmedo.
  \end{itemize}
  \vspace{0.2cm}
  \footnotesize Se dictan por ahora con el material de cátedra existente (\texttt{notebooks\_originales/10-15});
  quedan priorizados para la siguiente fase de transformación a OVA.
\end{frame}

% =================================================================
\section{Motivación}
% =================================================================
\begin{frame}{Lo que ya ha salido de este curso}
  \begin{itemize}
    \item Este no es el primer semestre en que un proyecto de Máquinas Térmicas se toma en serio.
    \item Varios equipos y trabajos de grado de semestres anteriores llegaron a resultados publicables,
      prototipos físicos y herramientas con relevancia de industria.
    \item Cuatro ejemplos reales, del propio archivo del curso:
  \end{itemize}
\end{frame}

\begin{frame}{Ejemplo 1 — de un proyecto a un artículo científico}
  \textbf{``One-Dimensional Numerical Modelling of a Supersonic RAMJET Engine with Parametrizable Geometry''}
  \small (S.\ Cadavid Franco, C.\ Bayona Roa)
  \begin{itemize}
    \item Modelo 1D en MATLAB de un motor RAMJET supersónico, geometría parametrizable, validado contra
      literatura (Dallos).
    \item Resultados: empuje neto \textbf{1719.4 N}, impulso específico \textbf{2094.6 s}, Mach de salida
      \textbf{2.033}, eficiencia global \textbf{30.3\%}.
    \item Incluye análisis de sensibilidad completo (Mach de vuelo, geometría del difusor, relación
      combustible-aire, recuperación de presión).
    \item Es exactamente el tipo de contenido del \textbf{Bloque 6} (turbinas de gas / propulsión) que aún
      no tiene su OVA — candidato natural para tu proyecto.
  \end{itemize}
\end{frame}

\begin{frame}{Ejemplo 2 — de un trabajo de grado a un prototipo físico}
  \textbf{TG253015 — ``Desarrollo tecnológico de un motor rotativo Wankel de 6 tiempos mediante la
  modelación termo-mecánica de su funcionamiento''}
  \begin{itemize}
    \item Motor Wankel de 6 tiempos con segunda combustión controlada, dos variantes geométricas
      (3 lóbulos/4 caras y 4 lóbulos/5 caras), modelado 100\% en Python.
    \item Tres configuraciones comparadas: válvula de expansión, sobrealimentado, inyección de agua.
    \item Resultado: la variante con válvula de expansión alcanza \textbf{68.3\%} de eficiencia térmica,
      frente a \textbf{30.9\%} del motor convencional de 4 tiempos.
    \item Video del concepto y anexos completos disponibles en el archivo del curso — un trabajo de grado
      que empezó como proyecto de Máquinas Térmicas.
  \end{itemize}
\end{frame}

\begin{frame}{Ejemplo 3 — proyecto con pertinencia de mercado nacional}
  \textbf{[241030] ``Simulation of Turboprop Systems using Simulink: Development and Validation of
  Computational Mechanical Model''}
  \begin{itemize}
    \item Herramienta computacional (Simulink + teoría BEMT) para simular un motor turbohélice bajo
      distintas condiciones de vuelo y combustibles.
    \item Validada contra datos reales de turbohélice (p.\,ej.\ familia PT6/ATR).
    \item Los propios autores lo plantean como soporte para \textbf{mantenimiento, manufactura y
      reconversión energética a nivel nacional} — exactamente la lógica de la \textbf{Idea 4} (modelo de
      negocio de pertinencia nacional).
  \end{itemize}
\end{frame}

\begin{frame}{Ejemplo 4 — los proyectos de curso también dejan huella}
  \footnotesize
  Entregas de semestres recientes, con autoría de sus propios estudiantes:
  \begin{itemize}
    \item \emph{Planta Termoeléctrica} — cogeneración, equipos hidráulicos y termodinámicos (J.A.\ Sánchez,
      J.D.\ Aparicio, J.D.\ Roa).
    \item \emph{Energía Nuclear} — adaptado como OVA 05b, módulo de frontera (A.G.\ Belalcázar, S.F.\ Triana,
      S.\ Lemus, J.C.\ Pineda).
    \item \emph{Modelamiento termo-fluido de un intercooler aire-aire} acoplado a un motor diésel (J.C.\ Muñoz,
      S.G.\ Pedraza).
    \item \emph{Disipación térmica para dispositivos electrónicos} (J.E.\ Dávila, J.D.\ Saavedra).
    \item \emph{Análisis de combustibles para la Fórmula 1} (A.G.\ Belalcázar, S.\ Libos Dallos).
    \item \emph{Intercooler en un motor Wankel 6T turboalimentado} — taller de tercer corte.
  \end{itemize}
  \vspace{0.2cm}
  Y el artículo sobre combustión de EFB/PMF de palma (Arias-Guerrero, Castillo-Hernández, Mayorga-Guzmán,
  Bayona-Roa) hoy es el \textbf{OVA 08b} de este mismo curso.
\end{frame}

% =================================================================
\section{Proyección}
% =================================================================
\begin{frame}{¿Hasta dónde puede llegar tu proyecto?}
  \centering
  \begin{tabular}{cccc}
    \textbf{Proyecto de curso} & $\rightarrow$ & \textbf{Semillero de Máquinas} & \\
    \small (maqueta + negocio) & & \small investigación aplicada &  \\[0.4cm]
  \end{tabular}

  \begin{tabular}{ccc}
    \textbf{Trabajo / Proyecto de grado} & $\rightarrow$ & \textbf{Artículo científico} \\
    \small ej.\ TG253015 (Wankel 6T) & & \small ej.\ RAMJET, EFB/PMF (OVA 08b) \\
  \end{tabular}
  \vspace{0.5cm}

  \raggedright
  \begin{itemize}
    \item El camino \textbf{ya está recorrido} por otros equipos de este curso — no es hipotético.
    \item Un buen proyecto de maqueta puede convertirse en línea de trabajo del \textbf{Semillero de
      Máquinas}, en \textbf{trabajo o proyecto de grado}, o incluso en \textbf{publicación} co-autoría con
      el profesor.
    \item La calidad técnica (memoria de cálculo, validación, trazabilidad CAD-fabricación) es lo que abre
      esa puerta.
  \end{itemize}
\end{frame}

% =================================================================
\section{Cronograma}
% =================================================================
\begin{frame}{Hitos del proyecto en el semestre}
  \begin{itemize}
    \item \textbf{Semana 1} — Entrega de OVA en Teams; conformación de equipos.
    \item \textbf{Semana 6} — Cierre Corte 1: kick-off formal del proyecto, memoria de cálculo preliminar.
    \item \textbf{Semana 10} — Memoria de cálculo consolidada + diseño CAD del motor.
    \item \textbf{Semana 12} — Cierre Corte 2: avance de fabricación en el CTAI.
    \item \textbf{Semana 16} — Cierre Corte 3: entrega de la maqueta ensamblada y sustentación del modelo
      de negocio.
  \end{itemize}
\end{frame}

% =================================================================
\section{Cierre}
% =================================================================
\begin{frame}{Lo que se espera de cada equipo}
  \begin{itemize}
    \item Preparar cada OVA \textbf{antes} de la sesión presencial.
    \item Participar activamente en la discusión invertida.
    \item Sostener el rigor técnico de la memoria de cálculo.
    \item Liderar el proyecto como un equipo real, con roles y cronograma.
    \item Defender la maqueta \emph{y} su modelo de negocio.
  \end{itemize}
  \vspace{0.3cm}
  Los ejemplos de hoy (RAMJET, Wankel 6T, Turboprop, EFB/PMF) empezaron como proyectos de este curso.
  \textbf{El próximo podría ser el tuyo.}
  \vspace{0.3cm}
  \centering \Large ¡Bienvenidos al curso!
\end{frame}

\end{document}
